\section{Preparing the crossed large-mass limit}
\label{sec:crosslimit}
The large-mass limit must specify whether an observer threshold moves
with the prepared state. We first obtain estimates independent of the
large mass, then isolate the remaining local Racah limit. The latter is
stated as a hypothesis, not inferred from a pointwise asymptotic formula.

\subsection{Two distinct preparation protocols}
Use exact observer energy $m(s)=\sqrt{s^2+1/4}$, its inverse
$s(m)=\sqrt{m^2-1/4}$, and $j(m)=m/s(m)$. Let $\pi_0(r)$ be a fixed
probability density on $r\ge0$. A translated normalized vacuum packet is
specified by
\begin{equation}
 \rhoP(d)|f_M(d)|^2\dd d=\pi_0(r)\dd r,
 \qquad m(d)=M+r.
 \label{eq:translated-energy-packet}
\end{equation}
The phase of $f_M$ is immaterial to these pure-right vacuum correlators.
The earlier flat clock variable is $p=2\pi r$; if its density is $P(p)$,
then $\pi_0(r)=2\pi P(2\pi r)$.
We distinguish two compressions applied at every insertion:
\begin{equation}
 Q^{\rm fixed}=\one_{H_R\ge m_*},\quad m_*>1/2;
 \qquad Q^{\rm moving}_M=\one_{H_R\ge M},\quad M\ge1.
 \label{eq:two-thresholds}
\end{equation}
The left observer is already in the support of the packet. Its matching
threshold projection acts trivially throughout a pure-right correlator.
In excess energy the first boundary retreats to $m_*-M$, whereas the
second remains at zero. The fixed-threshold problem retains our positive
$s_*$ assumption and does not assert the $s_*=0$ field extension.

The microcanonical preparation of~\cite[(6.29), (B.43)]{KL} has a window
on the external state, without a moving threshold at every intermediate
insertion. Its variables $s_i=\Lambda+q_i$ are representation labels;
our $M+r_i$ are exact energies. On bounded offsets, differences of the
two labels agree with energy differences up to $O(M^{-2})$.
The paper's fixed-time delta limit and its scrambling-time limit are
different distributional limits. Its two time origins are
$\log(2\Lambda)$ in (6.32) and $\log(2\Lambda/e)$ in (B.45).
Neither distinction is removed by a normalization convention for our
packet.

\subsection{Uniform tails without an absolute bound on the Racah kernel}
The one-step positive density in exact energy is
\begin{equation}
 K(m'|m)=j(m')V_\nu(s(m')|s(m))
                  e^{-\sigma^2(m'-m)^2}.
 \label{eq:energy-positive-kernel}
\end{equation}
Fix $\nu,\sigma>0$ and a common lower mass $m_*>1/2$ for all states
under consideration. Equation~\eqref{eq:gamma-envelope} and
$j(m')\le j(m_*)$ give a constant $C$ such that
\begin{equation}
 0\le K(m'|m)\le B(m'-m):=C e^{-\sigma^2(m'-m)^2}.
 \quad A:=\int_{\mathbb R}B(u)\dd u<\infty.
 \label{eq:uniform-energy-envelope}
\end{equation}
Additional intermediate projections only reduce positive norms.
For a vacuum fiber at energy $m_0$, the squared norm of either ordered
one-of-each state is the integral of
$K(m_1|m_0)K(m_2|m_1)$ and is at most $A^2$.
This uses the normalized ordered CG tree, not orthogonality of arbitrary
noncommuting projection histories.

\begin{proposition}[Mass-independent spectral-tail control]
In each ordered tree retain only
$|m_1-m_0|\le R$ and $|m_2-m_0|\le R$. Let $X_{M,R}$ pair the two
retained trees, and let $X_M$ be the uncut crossed pairing. Define
$b(R)=\int_{|u|>R}B(u)\dd u$. For either protocol in
\eqref{eq:two-thresholds}, any normalized packet, and all real times,
\begin{equation}
 |X_M-X_{M,R}|\le 2A\sqrt{2A\,b(R/2)}.
 \label{eq:uniform-cross-tail}
\end{equation}
Discarding additionally the external packet tail $r>L$ changes the
crossed pairing by at most
$A^2\int_L^\infty\pi_0(r)\dd r$.
\label{prop:uniform-cross-tail}
\end{proposition}
\begin{proof}
Write $u=m_1-m_0$, $v=m_2-m_1$. Leaving the retained path implies
$|u|>R/2$ or $|v|>R/2$. Ordered-channel orthogonality identifies its
squared norm with the positive excluded path integral. Thus it is at
most $2A b(R/2)$, uniformly in $m_0$ and time. Subtract the two pairings
and apply Cauchy--Schwarz, using norm at most $A$ for each tree.
The fiber bound $|X_{m_0}|\le A^2$ gives the external-tail assertion.
\end{proof}
The cutoff in this proposition is a numerical path restriction centered
on the external energy. It is not an additional fixed experimental mass
window. In particular, the old box $[0.2,3]$ cannot follow $M\to\infty$.
The constants are analytic finite constants; no certified numerical
value of $C$ is claimed.

\subsection{The unresolved compact distribution limit}
Let $\mathcal C_\nu(a,b,c;d)=X_H(a,b,c,d)/\rho_H(d)$ with exactly the
normalization in Section~\ref{sec:fourpoint}. In the variables
\begin{equation}
 d=s(M+r),\quad a=s(M+r+x),\quad
 b=s(M+r+y),\quad c=s(M+r+z),
\end{equation}
define the density relative to $\dd x\,\dd y\,\dd z$ by
\begin{equation}
 \widetilde{\mathcal C}_M(r;x,y,z)
 =j(M+r+x)j(M+r+y)j(M+r+z)\mathcal C_\nu(a,b,c;d).
 \label{eq:translated-cross-density}
\end{equation}
All masses in a compact set of offsets are positive for sufficiently
large $M$. The natural local hypothesis is
\begin{equation}
 \widetilde{\mathcal C}_M
 \longrightarrow v_\nu(x)v_\nu(z)\delta(y-x-z)
 \quad\text{in }\cD'(\mathbb R^4_{r,x,y,z}),
 \label{eq:local-racah-hypothesis}
\end{equation}
where
\begin{equation}
 v_\nu(u)=\frac{e^{-\pi u}}{16\pi^3}
 \left|\Gamma\left(\tfrac14+\tfrac i2(\nu+u)\right)
        \Gamma\left(\tfrac14+\tfrac i2(\nu-u)\right)\right|^2,
 \qquad k(u)=v_\nu(u)e^{-\sigma^2u^2}.
 \label{eq:limit-step-weight}
\end{equation}
The normalization is consistent with the assembled factors (B.41),
(B.42), (B.46) of~\cite{KL}: their conditional density on the delta
surface has coefficient
$e^{-\pi(x+z)}\Gamma_4(\nu,x)\Gamma_4(\nu,z)/(256\pi^6)$.
This coefficient check is not a proof of
\eqref{eq:local-racah-hypothesis}.

There is a specific unresolved analytic step. The exact ${}_4F_3(1)$
series has excess one. The ${}_2F_1(1)$ appearing in (B.33) has
excess $i(s_1-s_2+s_3-s_4)$, with zero real part. Its gamma expression
therefore involves continuation, not ordinary absolute summation at
one. Uniform control of the four combined branches near that excess
being zero is required. The fixed-label simplification alone does not
bound the remainder against a compact smooth test, especially near the
delta surface and coincident-label charts.

\begin{proposition}[Conditional reduction to a finite-energy problem]
If \eqref{eq:local-racah-hypothesis} holds, the exact crossed correlators
of both protocols converge for every fixed probability density $\pi_0$.
The convergence is uniform on compact sets of the four real times. For
the moving threshold the limit is
\begin{align}
 X_\infty(\mathbf t)
 &=\int_{\mathbb R^2}k(x)k(z)
      K_0((-x)_++(-z)_+)
      e^{ix(t_1-t_3)+iz(t_2-t_4)}\dd x\dd z,\notag\\
 K_0(q)&=\int_q^\infty\pi_0(r)\dd r,\qquad q\ge0.
 \label{eq:moving-cross-candidate}
\end{align}
For the fixed threshold, $K_0$ in this expression is replaced by one.
No rate or logarithmically growing time interval follows from this
proposition.
\label{prop:conditional-cross-limit}
\end{proposition}
\begin{proof}
First truncate external $r$ and path offsets by
Proposition~\ref{prop:uniform-cross-tail}. Approximate $\pi_0$ in $L^1$
by smooth compact densities; the fiber norm bound controls the error by
$A^2$ times that $L^1$ error. For the moving boundary, replace each
intermediate hard step by a smooth step differing only in an energy
layer of width $\eta$. The excluded norm squared is at most
$2\eta A\sup B$ per tree, because either the first or the final energy
is in the layer. Thus the pairing error is $O(\sqrt\eta)$ uniformly in
$M$. External layers are controlled by the mass of $\pi_0$ in them.
Smooth cutoffs on the retained offsets can similarly be chosen, or the
outer cutoff can be enlarged. We now have a compact smooth test for
\eqref{eq:local-racah-hypothesis}.

The exact Gaussian and time factors in these coordinates are
\begin{align}
 \mathcal Q&=x^2+(y-x)^2+(y-z)^2+z^2,\notag\\
 \mathcal T&=-xt_1+(x-y)t_2+(y-z)t_3+zt_4.
\end{align}
On $y=x+z$ they give $k(x)k(z)$ and the phase in
\eqref{eq:moving-cross-candidate}. The four surviving moving-threshold
conditions are $r,r+x,r+z,r+x+z\ge0$, equivalent to
$r\ge(-x)_++(-z)_+$. The fixed boundary eventually leaves any compact
set of offsets, and gives no such condition. Remove the smooth layers,
then the truncations, using the uniform bounds and their limiting
positive-weight counterparts.

Finally, differentiating a tree's time phase introduces one of its
energy increments. Equation~\eqref{eq:uniform-energy-envelope} has a
finite second increment moment, so the tree derivatives and hence the
four-point Lipschitz constants are bounded independently of $M$.
A finite time net upgrades pointwise convergence to uniform convergence
on compact time sets. This argument supplies no growing time net with
controlled convergence rate.
\end{proof}

\subsection{A boundary contribution in the comparison model}
Independently of the unproved hypothesis, define a flat-energy comparison
model with a clock on $L^2(\mathbb R_+)$ and two independent particle
species labelled by increments $x\in\mathbb R$. An emission has amplitude
$\sqrt{k(x)}$, changes clock energy from $r$ to $r+x$, and is projected
onto $r+x\ge0$. Negative particle labels here encode frequencies; they
are not an assumed physical Hamiltonian spectrum of de Sitter matter.
On the vacuum and one-particle sectors this construction is bounded
since $k\in L^1$. Smearing times multiply each emission by $e^{-ixt}$.
It is a comparison of ordered amplitudes, not a constructed full
de Sitter dual.

For equal species parameters the two ordered norms and their crossed
overlap at $(t,0,t,0)$ are, respectively,
\begin{align}
 N_\infty&=\int k(x)k(z)
       K_0(\max(0,-x,-x-z))\dd x\dd z,\notag\\
 X_\infty&=\int k(x)k(z)K_0((-x)_++(-z)_+)\dd x\dd z.
 \label{eq:comparison-norms}
\end{align}
These are direct formulas in this explicitly defined model.
The second formula is also the conditional prediction of
Proposition~\ref{prop:conditional-cross-limit} for the exact problem.
Both are independent of $t$ at this repeated-time configuration.
Subtracting the two ordering indicators gives the positive identity
\begin{equation}
 \mathcal D_{\rm edge}:=2(N_\infty-X_\infty)
 =2\int_0^\infty\pi_0(r)\dd r\int_0^\infty k(x)\dd x
                                  \int_r^{r+x}k(-u)\dd u>0.
 \label{eq:comparison-boundary-floor}
\end{equation}
Indeed the only unequal region is an energy gain $x>0$ followed by a
loss $u>r$, while $u\le r+x$ allows the final clock state. The reverse
order is removed at its first step. Since $k$ is strictly positive at
all finite real increments, the inner integrals are positive for every
finite $r$, proving strict positivity for any probability density.
This is the commutator norm squared in the comparison model; its
identification with a large-mass limit of the exact commutator remains
conditional on \eqref{eq:local-racah-hypothesis}.

If a fixed packet is shifted by an additional excess energy $h\to\infty$,
dominated convergence in \eqref{eq:comparison-boundary-floor} gives
$\mathcal D_{\rm edge}\to0$. Removing the clock boundary altogether
gives $N_\infty=X_\infty=(\int k)^2$. A nonzero repeated-time
commutator for a moving threshold can therefore include a static
projection contribution. It cannot by itself diagnose recoil growth or
a Lyapunov exponent. Establishing the local distribution limit and then
its time-dependent remainder is the next exact-model task.
